% Source: tex/chapters/ch06/04_ソフトウェアエンジニアリング運用の基礎.tex
\subsubsection{ソフトウェアエンジニアリング運用の定義}

ソフトウェアエンジニアリング運用とは、ソフトウェア製品、情報技術インフラ、システムソフトウェア、アプリケーションソフトウェアが、開発中、保守中、実運用中に適切に動作するように使われる知識、技能、プロセス、ツールの総体である。

運用担当技術者は、単にサーバを見張る人ではない。コンテナや仮想サーバの準備、配備、設定、支援を行う。必要な環境をオンデマンドで提供し、継続的インテグレーション、試験、配備、監視をサービスとして提供する。障害を診断し、問題と解決策を記録し、優先順位を付け、影響を評価する。さらに、セキュリティ、データ保護、フェイルオーバー、容量、ストレージ、データベース性能、技術仕様書の整備にも関わる。

\par\smallskip\noindent\refstepcounter{table}\label{tab:ch06-02}\textbf{表 \thetable: 第06章 ソフトウェアエンジニアリング運用 表02}\par\smallskip
\begin{longtable}{L{45mm}L{105mm}}
\toprule
役割 & 主な責任 \\
\midrule
運用担当技術者 & 運用サービスを設計し、配備、監視、容量管理、障害対応、セキュリティ、復旧を実行または自動化する。 \\
ソフトウェア技術者 & 運用担当技術者が提供する運用サービスを使い、自分のソフトウェアを配備、監視、改善する。 \\
プラットフォーム技術者 & 開発者が自律的に使える共通基盤、セルフサービス環境、配備基盤、監視基盤を作る。 \\
SRE技術者 & 可用性、性能、遅延、容量、緊急対応、変更管理を、測定と自動化に基づいて改善する。 \\
\bottomrule
\end{longtable}
